¿Qué es spoofing?

El spoofing es una técnica de suplantación de identidad digital maliciosa en la que el ciberdelincuente se hace pasar por una persona, dispositivo o entidad de confianza para engañar a los usuarios, sistemas o redes con el fin de obtener información confidencial, distribuir malware o facilitar ataques más avanzados~\cite{spoof}.

Existen distintos tipos de spoofing:
\begin{itemizeColor}
    \item \textbf{IP spoofing}: Falsificación de dirección IP.
    \item \textbf{Email spoofing}: Imitación de direcciones de correo electrónico.
    \item \textbf{DNS spoofing}: Redirección a sitios falsos.
    \item \textbf{ARP spoofing}: Manipulación de tablas de red.
\end{itemizeColor}

Para prevenir los ataques de spoofing es necesario implementar protocolos de autenticación, cifrado, filtros antiphishing y detección avanzada de anomalías, así como entrenar al personal para reconocer comportamientos sospechosos.
